\scanpage{108}
\mspage{99}
% archive sheet number: 108

Cas de ``\unsure{raccord}'', \deleted{quand $\unsure{\Phi|\varphi}$} En d'autres termes, on exclut le cas où $\Phi,\Psi$ seraient des parties \inserted{\unsure{à la fois}} ouvertes et fermées de $\Phi\cup\Psi$.

Ceci posé, voici la structure d'une préfigure quelconque, en termes des préfigures connexes.

\deleted{Si on a}
\inserted{Soit}
\[
D=\{\Phi_1\underset{\deleted{\mathrm{pos}}}{\overset{\unsure{!|'|1}}{\unsure{\ll|<}}}\Phi_2
\underset{\deleted{\mathrm{pos}}}{\overset{!}{\ll}}\cdots
\underset{\deleted{\mathrm{pos}}}{\overset{!}{\ll}}\Phi_n\}
\]
un drapeau \deleted{\unsure{pour}} de $\operatorname{Préfig}_{\mathrm{conn}}(\deleted{\unsure{\Phi}}\unsure{\mathscr{M}|\mathcal{M}})$
\[
D\in{\unsure{\operatorname{Drap}}}_{<^!}^{*}(\operatorname{Préfig}_{\mathrm{conn}}(\mathscr{M}))
\]
on lui associe la préfigure
\[
\Phi_D=\bigcup_{\Phi\in D}\Phi\ \deleted{\unsure{\amalg}}\ =\bigcup_{i\in[1,n]}\Phi_i.
\]
(réunion disjointe). On a fait ce qu'il fallait pour que les $\Phi_i$ soient les comp. connexes de $\Phi$. On \unsure{trouvera} que dans le cas général \deleted{\unsure{les}} les figures s'obtiennent comme ça, i.e.

\uline{Lemme} Soient $\Phi,\Psi$ préfigures connexes, Pour que $\Phi\cap\Psi=\varnothing$ et $\Phi\cup\Psi$ soit une préfigure, \deleted{\unsure{et on a|il suffit} $\Phi\underset{\mathrm{pos}}{<}\Psi$,} il \unsure{f et s|f. et s.} que $\Phi\underset{\mathrm{pos}}{<}\Psi$. Pour que $\Phi$ \deleted{et $\Psi$ soient}
